\section{Sharp boundaries of the proved comparisons}\label{sec:boundaries}
The following witnesses isolate distinct hypotheses. They are mathematical counterexamples to specific uncharged inferences, not a replacement of the main theorem by an impossibility statement.

\begin{proposition}[Absolute task error can be amplified by selection]\label{prop:absolute-sharp}
For every $b\ge1$ and $0\le e\le1/b$, there are positive prepared and actual joint laws with identical report marginals and one-call absolute squared-task defect at most $e$, whose optimal stopped risks at expected call budget $b$ differ by $be$, even with one output value. The difference occurs in either direction. A fixed relative task order is therefore strictly stronger than a small absolute task defect near a zero-loss component.
\end{proposition}
\begin{proof}
Draw an independent uniform report $x\in[0,1]$ at every call and let $D$ have mass $1/b$. All audits have conditional mean zero. Off $D$, both laws have a fresh fair-sign audit of magnitude one. On $D$, the comparison has a fair sign of magnitude $\sqrt{be}$ and the actual law has audit zero. For every $q\in[0,1]$ and finite-image readout $r(x)$, their unnormalized squared losses differ by $be\int_Dq(x)dx\le e$. The comparison residual variance is everywhere at least $be$, so every committing strategy has risk at least $be$; selecting $D$ and outputting zero attains it with expected cost $b$. The actual law attains zero by the same procedure. Reversing the laws proves the other direction. A single event-stopped label can reject at itself and commit with a final event; a state-stopped realization uses an additional live label. On $D$ a task zero under one law is positive under the other, so no two-sided relative order with a positive lower factor can hold.
\end{proof}

A report-only comparison also fails: a constant mark can have the same marginal under two laws while an additional visible bit reveals a fair audit sign. A two-output decoder then has zero loss, although a profile formed from the constant mark alone has residual variance one. The full-report conditional moments in \eqref{eq:full-moments} exclude this omission.


There is also an all-program witness for the necessity of more than report balance. Take the three admitted models $\theta\in\{-1,0,1\}$, let every precommit report be the same constant, and let the unseen audit at commitment be $Z=\theta$. There is one legal productive command, costing one, and no command changes $\theta$. The actual comparison member $\theta=0$ has $\Q_m(b)=0$. Nevertheless the distribution of every committed readout $R$ is identical under $\theta=\pm1$, because all precommit laws, including stopping and coins, are identical. The average of their risks is $1+\E R^2\ge1$, and output zero attains worst-case risk one. This holds for every almost-surely terminating policy, at every feasible budget and every dictionary size. Report masses match exactly; the missing first task moment is the obstruction. It is not repaired by choosing a different causal rejection rule.

A live-state count alone cannot justify a terminal-width subtraction. In the two-port interface, use labels $0,1$ with live commands gate and productive respectively, and decoder values $-1,+1$. Gate failure continues at $0$, success at $1$; a rejected productive report returns to $0$. A productive commitment emits a final event with label $0$ or $1$ according to its chosen output. Two live and two output roles have used only two labels. The state-stopped model cannot share those roles because its halt map is intrinsic. This is the exact semantic repair retained from the earlier work.

Finally, bounded-law discrepancy is not a cost certificate. A source procedure can agree with a target terminal output with probability one while performing an independent preliminary waiting time of arbitrarily large finite mean, or of infinite mean. Terminal total variation is then zero and expected cost is uncontrolled. Our simulator definition includes the scored stopping information when its law is compared, and still requires a separate expected-cost certificate because time is unbounded. Likewise, neither fixed-time equivalence nor finite-dimensional positive matrix identities imply large-deviation equivalence, unrestricted optimal-control equivalence, or closure of unbounded operators.

\subsection*{What is and is not generalized}
The central theorem generalizes exact and absolute-calibrated causal profile transfer to a common affine task order and derives its all-scale relative and stratified consequences. The positive-matrix characterization concerns localized task order; the quotient construction concerns compact-exhausted continuous-response experiments; the pilot concerns a fixed independent command sampler. These results do not establish necessity of a common acquisition program for every causal experiment, nor the minimal possible simulator state under all implementations, nor matching lower bounds for program length, scratch and arithmetic time. Those stronger mathematical questions are distinct from the fully stated sufficient-class theorem proved here. The general theorem does not depend on a conclusion taken from either realization. The retained nonintervening sensing, occupation-flow, exact role and absolute-defect results remain substantive independent parts of the complete companion.
